\documentclass[../main.tex]{subfiles}
\begin{document}

%========================================================================================
% FILENAME: 08_seigniorage.tex (v1.0.0)
%
% §8 Seigniorage — the bootstrapping bound under immediate burning, and the negative
% lifetime result.
%
% Proof budget: two body proofs (bootstrapping-bound; no-finite-lifetime-envelope).
% No \interfaceref.  Silences applied; cycle only; "non-operator agents".  First fee-flow
% mention uses \feeflow.
%
% Change history lives in version control (see the versioning policy in main.tex).
%========================================================================================

\section{Seigniorage}
\label{sec:attestation:seigniorage}

The operator's only inflow is the minting fee, a per-cycle flow of $\feeflow$ during the
bootstrapping phase. This chapter bounds the cumulative attestation the operator can extract
from it during bootstrapping under immediate burning, and shows that no such bound exists
over the receipt's lifetime.

\begin{theorem}[Seigniorage Bound --- Bootstrapping Phase]
    \label{thm:seigniorage:bootstrapping-bound}
    Let $D$ denote cumulative reserve units deposited by non-operator agents during the
    bootstrapping phase, and suppose every fee-derived \liveclass{} unit is burned
    \emph{immediately, at the floor prevailing when it was minted}. Then operator cumulative
    attestation is
    \[
        \attest{\opidx} = \feerate\,\livesplit\,D.
    \]
    The immediacy hypothesis is not decorative: it is exactly what fixes the valuation
    floor. Without it no finite envelope of this form exists
    (\zcref{thm:seigniorage:no-finite-lifetime-envelope}).
\end{theorem}
\begin{proof}
    Fee units per minting batch $\Delta\reservepool{}$ are
    $\feerate\livesplit\Delta\reservepool{}/\ratefloor{}^{-}$
    (\zcref{alg:operations:cycle-processing}); burned at the minting floor, the attestation
    is $\feerate\livesplit\Delta\reservepool{}$. Summing over cycles gives
    $\feerate\livesplit D$.
\end{proof}

\begin{theorem}[No Finite Lifetime Envelope]
    \label{thm:seigniorage:no-finite-lifetime-envelope}
    There is no finite bound on $\attest{\opidx}$ over the receipt lifetime depending only
    on $\feerate$, $\livesplit$, and $D$.
\end{theorem}
\begin{proof}
    Maturity conversion sets $\tlocksupply{} = 0$
    (\zcref{def:maturity:conversion}), after which $\ratefloor{}$ is unbounded above
    (\zcref{prop:containment:time-locked-standing-bound}). An operator that defers burning
    until other holders have burned first therefore settles at an arbitrarily high floor.

    Concretely at $\livesplit = \feerate = 1/2$, $\genesisreserve = 1000$, $D = 1000$: after
    conversion $\reservepool{} = 2000$ and $\totsupply{} = 2000$, of which the operator holds
    $500$ fee-derived units. Other holders burning their $1500$ leaves
    $\ratefloor{} = 4$; the operator burning $250$ of its holding already credits $1000$,
    and burning a further $249$ at $\ratefloor{} = 8$ credits $1992$ more. More generally,
    burning supply from $\totsupply{}\sp{-}$ down to $\varepsilon$ credits
    $\reservepool{}\ln(\totsupply{}\sp{-}/\varepsilon)$, which diverges as
    $\varepsilon \to 0\sp{+}$.
\end{proof}

\begin{remark}[What Replaces the Envelope]
    \label{rem:seigniorage:withdrawn-envelope}
    Earlier revisions of this document carried a finite lifetime envelope
    $\attest{\opidx} \leq \feerate D/(1-\livesplit)$. It was false: its derivation valued
    post-maturity burns at the pre-deposit floor factor
    (\zcref{cor:containment:pre-deposit-floor-factor}) while the same document states that
    $\ratefloor{}$ is unbounded above post-maturity. Its label is withdrawn rather than
    reused, since a stable anchor must not change meaning
    \atthislayer{}. What survives is
    \zcref{thm:seigniorage:bootstrapping-bound} under its immediacy hypothesis, the
    pre-maturity capacity bound
    (\zcref{thm:containment:pre-maturity-deposit-capacity}), and the negative result above.
\end{remark}

\stackmarginnote{Note:}{Fee-channel: self-resolving by
\zcref{prop:discipline:anti-correlation}.}%
No-recycling-advantage (\zcref{prop:costliness:no-recycling-advantage}) extends to the
time-locked class post-maturity: once converted, a time-locked-derived \liveclass{} unit
behaves identically to any other \liveclass{} unit, including the $W$-conserving recycling
bound. During bootstrapping the time-locked class cannot be burned, so no recycling through
it is possible.

%========================================================================================
% END OF FILE: 08_seigniorage.tex (v1.0.0)
%========================================================================================

\end{document}
